\subsection{THE AXIOM OF THE SET OF SUBSETS}

\[
(\forall \mathbf {X}) \operatorname{Coll} _ {\mathbf {Y}} (\mathbf {Y} \subset \mathbf {X}).
\]

This axiom means that for every set X there exists a set whose elements are all the subsets of X, namely the set \(\mathcal{E}_{\mathbf{Y}}(\mathbf{Y}\subset\mathbf{X})\) (§1, no. 4); this set is denoted by \(\mathfrak{P}(\mathbf{X})\) and is called the set of subsets of X. Clearly, if \(X\subset X'\) , we have \(\mathfrak{P}(\mathbf{X})\subset\mathfrak{P}(\mathbf{X}')\) .

Let A, B be two sets and let \(\Gamma\) be a correspondence between A and B. The function \(X \to \Gamma\langle X\rangle\) ( \(X \subset \mathfrak{P}(A)\) , \(\Gamma\langle X\rangle \in \mathfrak{P}(B)\) ) is called the canonical extension (or simply the extension) of \(\Gamma\) to sets of subsets and is denoted by \(\hat{\Gamma}\) ; it is a mapping of \(\mathfrak{P}(A)\) into \(\mathfrak{P}(B)\) . If \(\Gamma'\) is a correspondence between B and a set C, the formula \((\Gamma' \circ \Gamma)\langle X\rangle = \Gamma'\langle\Gamma\langle X\rangle\rangle\) shows that the extension of \(\Gamma' \circ \Gamma\) to sets of subsets is the mapping \(\hat{\Gamma}' \circ \hat{\Gamma}\) .

PROPOSITION 1. (1) If \(f\) is a surjection of a set \(\mathbf{E}\) onto a set \(\mathbf{F}\) , the canonical extension \(\hat{f}\) of \(f\) is a surjection of \(\mathfrak{P}(\mathbf{E})\) onto \(\mathfrak{P}(\mathbf{F})\) .

\begin{enumerate}
\def\labelenumi{(\arabic{enumi})}
\setcounter{enumi}{1}
\item
  If \(f\) is an injection of \(\mathbf{E}\) into \(\mathbf{F}\) , the canonical extension \(\hat{f}\) of \(f\) is an injection of \(\mathfrak{P}(\mathbf{E})\) into \(\mathfrak{P}(\mathbf{F})\) .
\item
  If \(s\) is a section of \(f\) , then \(f \circ s\) is the identity mapping of \(\mathbf{F}\) , hence \(\hat{f} \circ \hat{s}\) is the identity mapping of \(\mathfrak{P}(\mathbf{F})\) ; therefore \(\hat{f}\) is surjective and \(\hat{s}\) is a section of \(\hat{f}\) (§3, no. 8).
\item
  The proposition is obvious if \(\mathbf{E} = \varnothing\) , because then \(\mathfrak{P}(\mathbf{E}) = \{\varnothing\}\) . If \(\mathbf{E} \neq \varnothing\) and if \(r\) is a retraction of \(f\) , then \(r \circ f\) is the identity mapping of \(\mathbf{E}\) , so that \(\hat{r} \circ \hat{f}\) is the identity mapping of \(\mathfrak{P}(\mathbf{E})\) ; therefore \(\hat{f}\) is injective, and \(\hat{r}\) is a retraction of \(\hat{f}\) (§3, no. 8).
\end{enumerate}

